\chapter{Cauchy 数列与收敛准则}
\label{chap:cauchy-criterion}

\begin{chapterguide}
极限定义以未知的候选 \(a\) 为参照。Cauchy 条件只比较尾部各项，因而可以在尚不知道极限时检验数列是否具有收敛所需的内部稳定性。收敛必然推出 Cauchy 性；反向成立则依赖值域完备。本章在实数中证明充要准则，并用有理逼近 \(\sqrt2\) 说明不完备数系中的失败。
\end{chapterguide}

\section{“项与项接近”代替“项与极限接近”}

若 \(a_n\to a\)，则对充分大的 \(m,n\)，
\[
 \abs{a_m-a_n}
 \le\abs{a_m-a}+\abs{a_n-a}
\]
很小。这表明收敛列的尾部各项彼此接近。反过来，若所有后项彼此接近，就得到一个不出现未知极限的判别条件。

“相邻项接近”仍不够。调和部分和满足
\[
 H_{n+1}-H_n=\frac1{n+1}\to0,
\]
但 \(H_n\to+\infty\)。Cauchy 条件要求任意两个充分靠后的项接近，而非只控制相邻项。

\begin{example}
若存在非负数列 \(r_k\) 使
\[
 \abs{a_{k+1}-a_k}\le r_k,\qquad
 \sup_{m>n}\sum_{k=n}^{m-1}r_k\longrightarrow0,
\]
则
\[
 \abs{a_m-a_n}\le\sum_{k=n}^{m-1}r_k
\]
说明 \((a_n)\) 具有 Cauchy 性。几何误差 \(r_k=Cq^k\) 是典型情形。
\end{example}

\section{Cauchy 数列的定义与量词结构}

\begin{definition}
实数列 \((a_n)\) 称为 \term{Cauchy 数列}，如果
\[
 \forall\varepsilon>0\ \exists N\in\N\
 \forall m,n\ge N,\qquad
 \abs{a_m-a_n}<\varepsilon.
\]
\end{definition}

同一个 \(N\) 必须同时控制两个独立指标 \(m,n\)。把条件写成
\[
 \forall\varepsilon>0\ \forall m,n\ \exists N\le\min\{m,n\}
\]
会失去尾部统一性。严格否定为
\[
 \exists\varepsilon_0>0\ \forall N\in\N\
 \exists m,n\ge N,\qquad
 \abs{a_m-a_n}\ge\varepsilon_0.
\]

\begin{proposition}[尾部直径刻画]
令
\[
 d_N=\sup\{\abs{a_m-a_n}:m,n\ge N\}\in[0,+\infty].
\]
则 \((a_n)\) 是 Cauchy 列，当且仅当 \(d_N\to0\)。
\end{proposition}

\begin{proof}
若为 Cauchy 列，对给定 \(\varepsilon\) 可使尾部任意两项距离小于 \(\varepsilon\)，故 \(d_N\le\varepsilon\)；用 \(\varepsilon/2\) 可取得严格误差。反之若 \(d_N\to0\)，取 \(N\) 使 \(d_N<\varepsilon\)，定义立即成立。
\end{proof}

对有界列，尾部直径等于
\[
 d_N=\sup_{k\ge N}a_k-\inf_{k\ge N}a_k,
\]
所以\chapterref{chap:limsup-liminf}的尾部振幅判据也是 Cauchy 判据。

\section{收敛数列必为 Cauchy 数列}

\begin{theorem}
若实数列 \(a_n\to a\)，则 \((a_n)\) 是 Cauchy 列。
\end{theorem}

\begin{proof}
给定 \(\varepsilon>0\)，取 \(N\) 使
\[
 n\ge N\quad\Longrightarrow\quad
 \abs{a_n-a}<\frac\varepsilon2.
\]
于是对任意 \(m,n\ge N\)，
\[
 \abs{a_m-a_n}
 \le\abs{a_m-a}+\abs{a_n-a}<\varepsilon.
\]
\end{proof}

证明只使用三角不等式，不使用实数完备性。因此它在任意度量空间中都成立。误差 \(\varepsilon/2\) 是把总距离预算分给两项与极限之间的距离。

\begin{corollary}
Cauchy 列的任意子列仍是 Cauchy 列。若一个 Cauchy 列有收敛子列，则原列收敛到同一极限。
\end{corollary}

\begin{proof}
子列继承性由其指标最终进入原列尾部。若 \(a_{n_k}\to a\)，给定 \(\varepsilon>0\)，先用 Cauchy 性使原列尾部任意两项距离小于 \(\varepsilon/2\)，再选同一尾部中的子列项使其与 \(a\) 的距离小于 \(\varepsilon/2\)。三角不等式给出原列整个尾部与 \(a\) 的距离小于 \(\varepsilon\)。
\end{proof}

\section{Cauchy 数列必有界}

\begin{theorem}
每个 Cauchy 实数列都有界。
\end{theorem}

\begin{proof}
在定义中取 \(\varepsilon=1\)，得到 \(N\)，使 \(m,n\ge N\) 时
\(\abs{a_m-a_n}<1\)。固定 \(n=N\)，则对 \(m\ge N\)，
\[
 \abs{a_m}\le\abs{a_N}+1.
\]
再把前 \(N-1\) 项的绝对值并入有限最大值，便得到全列统一界。
\end{proof}

逆命题不成立，\((-1)^n\) 是有界非 Cauchy 列。Cauchy 性还比“相邻差趋零”强，因为它控制任意长度的尾部累积误差。

\begin{proposition}
若 \((a_n)\)、\((b_n)\) 为 Cauchy 列，则它们的和、差、积仍为 Cauchy 列；若存在 \(\delta>0\) 使 \(\abs{b_n}\ge\delta\) 最终成立，则 \(a_n/b_n\) 也是 Cauchy 列。
\end{proposition}

\begin{proof}
和、差由三角不等式。乘积使用 Cauchy 列有界性和分解
\[
 a_mb_m-a_nb_n=a_m(b_m-b_n)+b_n(a_m-a_n).
\]
商先证明倒数列：
\[
 \abs{\frac1{b_m}-\frac1{b_n}}
 \le\delta^{-2}\abs{b_m-b_n},
\]
再用乘积结论。
\end{proof}

\section{实数中的 Cauchy 收敛准则}

\begin{theorem}[Cauchy 收敛准则]
实数列收敛，当且仅当它是 Cauchy 数列。
\end{theorem}

\begin{proof}
正向已在第三节证明。反向设 \((a_n)\) 为 Cauchy 列。它有界，所以由 Bolzano--Weierstrass 定理存在收敛子列 \(a_{n_k}\to a\)。第三节推论说明具有收敛子列的 Cauchy 列整体收敛到 \(a\)。
\end{proof}

反向证明中，Cauchy 性提供全列尾部的内部控制，有界性与 Bolzano--Weierstrass 定理产生一个候选极限，最后再把子列极限推广到原列。真正依赖实数完备性的环节是有界列能够抽取收敛子列。

\begin{theorem}[快速子列证明]
也可以不直接调用 Bolzano--Weierstrass 定理：从 Cauchy 列抽取子列 \(a_{n_k}\)，使
\[
 \abs{a_{n_{k+1}}-a_{n_k}}<2^{-k}.
\]
则该子列收敛，进而原列收敛。
\end{theorem}

\begin{proof}
递归选择尾部起点即可得到该子列。对 \(m>k\)，
\[
 \abs{a_{n_m}-a_{n_k}}
 \le\sum_{j=k}^{m-1}2^{-j}<2^{1-k}.
\]
可用嵌套闭区间或单调有界的绝对收敛部分和证明该快速子列存在极限；随后使用第三节推论。该路线仍在“快速子列有极限”处使用实数完备性。
\end{proof}

\section{有理数中准则失败的经典例子}

\chapterref{chap:rational-defects}构造了有理数列
\[
 a_n=\frac{k_n}{10^n},\qquad
 a_n^2<2\le(a_n+10^{-n})^2,
\]
并证明
\[
 0\le a_m-a_n<10^{-n}\qquad(m\ge n).
\]
因此它是 \(\Q\) 中的 Cauchy 列，却没有有理极限。

\begin{theorem}
“每个 Cauchy 列收敛”不是有序域公理的结论，也不是 Archimedes 性质的结论。
\end{theorem}

\begin{proof}
\(\Q\) 是 Archimedes 有序域。上述有理 Cauchy 列若在 \(\Q\) 中收敛到 \(q\)，平方极限定理给出 \(q^2=2\)，与 \(\sqrt2\notin\Q\) 矛盾。
\end{proof}

同一数列在 \(\R\) 中收敛到 \(\sqrt2\)。Cauchy 性只使用项之间的距离，所以在 \(\Q\) 与 \(\R\) 中完全相同；“极限属于哪个空间”则取决于值域是否包含缺失边界。

\section{Cauchy 准则与空间完备性的预告}

\begin{definition}
若一个度量空间中的每个 Cauchy 列都收敛到该空间内的点，则称该空间\term{完备}。
\end{definition}

这里提前使用最少的度量语言：度量 \(d(x,y)\) 满足非负性、分离性、对称性和三角不等式；Cauchy 条件把
\(\abs{a_m-a_n}\) 换成 \(d(x_m,x_n)\)。实数完备，有理数不完备。闭区间作为实数中的子空间是完备的，而开区间通常不完备，例如 \(1/n\) 在 \((0,1)\) 中是 Cauchy 列但极限 \(0\) 不在该空间。

\begin{proposition}
完备度量空间的闭子集仍完备。反之，一个度量空间的完备子集必为闭集。
\end{proposition}

\begin{proof}
设 \(F\) 为闭子集，\(F\) 中 Cauchy 列在母空间中收敛到某点 \(x\)；闭性保证 \(x\in F\)。反之，若 \(F\) 完备而 \(x\) 是 \(F\) 中序列的母空间极限，则该序列为 Cauchy 列，在 \(F\) 中还收敛到某点 \(y\)。极限唯一性给出 \(x=y\in F\)。
\end{proof}

一般空间的完备化将在度量空间编系统构造。\chapterref{chap:cauchy-construction}以有理 Cauchy 列等价类构造实数，已经给出这一过程的基本模型。

\section*{本章小结}
\addcontentsline{toc}{section}{本章小结}

Cauchy 条件用尾部两项之间的统一距离替代未知极限。收敛列必为 Cauchy 列，Cauchy 列必有界。在实数中，有界性产生收敛子列，Cauchy 性再把子列极限推广到全列；这个反向蕴含刻画了完备性。有理逼近 \(\sqrt2\) 说明，内部稳定性并不保证极限仍属于不完备子空间。

\section*{习题}
\addcontentsline{toc}{section}{习题}

\noindent\textbf{配置说明：}本章按I级枢纽章配置，共24道编号题，含45个独立训练单元，建议总用时6至9小时。\exerciserange{13.1}{13.17}为核心必做范围。

\subsection*{A组　定义与基本性质}
\begin{enumerate}[label=\textbf{13.\arabic*.}]
 \exerciseitem{13.1} \mustdo 写出 Cauchy 条件的严格否定，并用于证明 \((-1)^n\) 不是 Cauchy 列。
 \exerciseitem{13.2} \mustdo 证明收敛列为 Cauchy 列，标明 \(\varepsilon/2\) 的作用。
 \exerciseitem{13.3} \mustdo 证明 Cauchy 列有界，完整处理有限首项。
 \exerciseitem{13.4} \mustdo 证明 Cauchy 列的任意子列仍为 Cauchy 列。
 \exerciseitem{13.5} \mustdo 证明 Cauchy 列若有收敛子列，则原列收敛到同一极限。
 \exerciseitem{13.6} \mustdo 证明尾部直径趋零与 Cauchy 条件等价。
 \exerciseitem{13.7} \mustdo 给出相邻差趋零但非 Cauchy 的数列，并从定义验证失败。
\end{enumerate}

\subsection*{B组　运算、判据与不完备性}
\begin{enumerate}[label=\textbf{13.\arabic*.},resume]
 \exerciseitem{13.8} \mustdo 证明两个 Cauchy 列的和与积仍为 Cauchy 列。
 \exerciseitem{13.9} \mustdo 若 \(a_n\) 是 Cauchy 列且最终 \(\abs{a_n}\ge\delta>0\)，证明 \(1/a_n\) 是 Cauchy 列。
 \exerciseitem{13.10} \mustdo 证明若 \(\abs{a_{n+1}-a_n}\le2^{-n}\)，则 \((a_n)\) 为 Cauchy 列。
 \exerciseitem{13.11} \mustdo 推广上一题：若非负级数尾和满足
 \[
 \sup_{m>n}\sum_{k=n}^{m-1}r_k\to0
 \]
 且 \(\abs{a_{k+1}-a_k}\le r_k\)，证明 \(a_n\) 为 Cauchy 列。
 \exerciseitem{13.12} \mustdo 使用 Bolzano--Weierstrass 定理证明实数 Cauchy 准则。
 \exerciseitem{13.13} \mustdo 对\chapterref{chap:rational-defects}十进制逼近 \(\sqrt2\) 的有理列，直接验证 Cauchy 性并证明无有理极限。
 \exerciseitem{13.14} \mustdo 构造 \((0,1)\) 中没有区间内极限的 Cauchy 列。
 \exerciseitem{13.15} \mustdo 证明实数中的闭区间完备，开区间一般不完备。
 \exerciseitem{13.16} \mustdo 判断“每个有界数列都是 Cauchy 列”“每个有收敛子列的数列都是 Cauchy 列”的真伪。
 \exerciseitem{13.17} \mustdo 证明 Cauchy 列的上下极限相等，并由此再证 Cauchy 收敛准则。
\end{enumerate}

\subsection*{C组　结构与项目}
\begin{enumerate}[label=\textbf{13.\arabic*.},resume]
 \exerciseitem{13.18} \projectdo\ 【前置：\chapterref{chap:equivalent-completeness}完备性；预计用时：75分钟】从 Cauchy 列抽取满足
 \(\abs{a_{n_{k+1}}-a_{n_k}}<2^{-k}\) 的快速子列，并不经过 Bolzano--Weierstrass 定理完成收敛证明。
 \exerciseitem{13.19} \mustdo 证明完备度量空间的闭子集完备，完备子集在母空间中闭。
 \exerciseitem{13.20} \optionaldo 给出一个数列，它在通常距离下是 Cauchy 列；把距离改为离散距离后不再是 Cauchy 列。
 \exerciseitem{13.21} \optionaldo 若两个实数列满足 \(\abs{a_n-b_n}\to0\)，证明其中一个为 Cauchy 列当且仅当另一个为 Cauchy 列。
 \exerciseitem{13.22} \optionaldo 证明 Cauchy 列的有限线性组合仍为 Cauchy 列，并讨论无穷线性组合为何需要额外一致控制。
 \exerciseitem{13.23} \opendo 比较“序完备”“Cauchy 完备”和“区间套完备”在 Archimedes 有序域与一般度量空间中的关系。
 \exerciseitem{13.24} \projectdo\ 【前置：\chapterref{chap:cauchy-construction}；预计用时：90分钟】概括 \(\Q\) 的 Cauchy 完备化：等价关系、常值嵌入、稠密性和完备性的四个证明模块。
\end{enumerate}

\section*{第二编综合作业}
\addcontentsline{toc}{section}{第二编综合作业}

\noindent 本作业串联第7--13章，共12道题、约52个独立小问；建议用时8至12小时。
\begin{enumerate}[label=\textbf{II.\arabic*.}]
 \exerciseitem{II.1} \mustdo 对命题“\(a_n\to a\)”依次完成：写出定义、写出否定、写出邻域形式、证明有限项修改不影响结论。
 \exerciseitem{II.2} \mustdo 设 \(a_n=(n+\alpha)/(n+\beta)\)。讨论分母最终非零，给出极限及关于参数 \((\alpha,\beta)\) 在有界矩形上的一致误差估计。
 \exerciseitem{II.3} \mustdo 设 \(a_n,b_n\) 收敛。证明线性组合、乘积、商、绝对值和最大值的极限定理，并为每个必要条件给出一个删去后失败的例子。
 \exerciseitem{II.4} \mustdo 分析递推 \(a_{n+1}=(a_n+3/a_n)/2\)、\(a_1>0\)：不变区间、单调性、有界性、极限与误差机制。
 \exerciseitem{II.5} \mustdo 对有界数列完成三种抽取：收敛子列、单调子列、趋于上极限的子列；比较三种构造的输入与输出。
 \exerciseitem{II.6} \mustdo 构造一个部分极限集合恰为 \(\{-2,0,1\}\) 的数列，计算上下极限、尾部振幅，并判断是否 Cauchy。
 \exerciseitem{II.7} \mustdo 证明下列条件等价：\(a_n\to L\)；每个子列都有进一步子列趋于 \(L\)；\(\liminf a_n=\limsup a_n=L\)；尾部振幅趋零。
 \exerciseitem{II.8} \mustdo 对 \(\sum_{k=1}^n2^{-k}u_k\)（\(\abs{u_k}\le1\)）证明 Cauchy 性、统一有界性及尾部误差，并讨论 \(u_k\) 不收敛是否妨碍部分和收敛。
 \exerciseitem{II.9} \projectdo\ 【预计用时：90分钟】从确界原理出发重建链条：单调有界收敛、区间套、Bolzano--Weierstrass、Cauchy 准则。
 \exerciseitem{II.10} \mustdo 比较 \(\Q\) 与 \(\R\) 中同一个有理 Cauchy 列：Cauchy 性、可能极限、值域闭性与完备性的差异。
 \exerciseitem{II.11} \optionaldo 构造有界列 \(a_n,b_n\)，使二者各有两个部分极限，而 \(a_n+b_n\) 分别出现一个、两个和三个部分极限的三组例子。
 \exerciseitem{II.12} \opendo 以“尾部”为主线总结第二编：最终性质、完整尾部、尾部值集、尾部确界、尾部直径分别解决什么问题，哪些信息彼此可恢复。
\end{enumerate}

\section*{部分习题提示}
\begin{itemize}
 \item \exerciseref{13.7}：使用调和部分和。
 \item \exerciseref{13.17}：对 Cauchy 列，尾部上、下确界之差不超过尾部直径。
 \item \exerciseref{13.20}：离散距离为 \(d(x,y)=1\)（\(x\ne y\)）。
 \item 综合题II.4：这是 \(\sqrt3\) 的 Newton 迭代。
 \item 综合题II.8：任意尾和绝对值不超过几何尾和。
\end{itemize}
